\documentclass[11pt]{amsart}
\usepackage{amssymb}
\usepackage{url}
\usepackage[margin=1.2in]{geometry}

\newtheorem{theorem}{Theorem}
\newtheorem{proposition}[theorem]{Proposition}
\newtheorem{lemma}[theorem]{Lemma}
\newtheorem{corollary}[theorem]{Corollary}
\theoremstyle{remark}
\newtheorem{remark}[theorem]{Remark}

\title{Rank-one Horrocks over a local ring by determinants}
\author{Claude (Anthropic, Opus 5.5)}
\date{}

\begin{document}

\begin{abstract}
Let $A$ be a local ring and $L$ an ideal of $A[X]$ containing a monic
polynomial, such that the inclusion $L\subseteq A[X]$ has an $A$-linear
retraction. We give a short proof that $L$ is generated by a monic
polynomial: $A[X]/L$ is free over $A$, the characteristic polynomial $G$ of
multiplication by $X$ lies in $L$, and a surjection between free modules of
the same finite rank is an isomorphism, so $L=(G)$. Each step is
constructive. The companion note~\cite{Div} proves the same theorem by
division, without freeness or determinants.
\end{abstract}

\maketitle

\section{Statement}

All rings are commutative. $A$ is local: if $a+b$ is a unit, then $a$ or $b$
is a unit. We do not require $0\neq1$. For a monic $F\in A[X]$ of degree $d$,
every $p$ has a unique decomposition $p=Fq+r$ with $\deg r<d$, and $q$ and $r$
are $A$-linear in $p$.

\begin{theorem}\label{thm:main}
Let $L$ be an ideal of $A[X]$ containing a monic polynomial $F$, and let
$\rho:A[X]\to L$ be $A$-linear with $\rho(p)=p$ for $p\in L$. Then $L$ is
generated by a monic polynomial, and this polynomial is unique.
\end{theorem}

Uniqueness holds over any ring: two monic polynomials dividing each other are
equal~\cite[Lemma~8]{Div}. The proof of existence has three steps:
$A[X]/L$ is free of some rank $e$ (Section~\ref{sec:free}); the
characteristic polynomial $G$ of multiplication by $X$ on it is monic of
degree $e$ and lies in $L$ (Section~\ref{sec:cayley}); and every element of
$L$ of degree $<e$ is $0$, so $G$ divides $L$ (Section~\ref{sec:cramer}).

\section{Freeness}\label{sec:free}

Let $V$ be an $A$-module and $N\subseteq V$ a submodule; write $p\equiv p'$
for $p-p'\in N$. A \emph{presentation} of $V/N$ by $m$ generators consists of
$g_1,\dots,g_m\in V$ and $A$-linear forms $\kappa_1,\dots,\kappa_m$ on $V$ with
\[
p\equiv\textstyle\sum_t\kappa_t(p)\,g_t\quad\text{for all }p\in V,
\qquad \kappa_t(p)=0\quad\text{for }p\in N.
\]
It is a \emph{basis} if moreover $\kappa_i(g_j)=\delta_{ij}$. Then
$\sum_tc_tg_t\in N$ forces $c=0$, by applying $\kappa$.

The forms $\kappa$ say that $V/N$ is projective: they define $V/N\to A^m$,
and $p\mapsto\sum_t\kappa_t(p)g_t$ splits it. The next lemma is the local
freeness lemma (see~\cite{LQ}), in this language.

\begin{lemma}\label{lem:free}
If $A$ is local, then $V/N$ has a basis.
\end{lemma}

\begin{proof}
Induction on $m$. For $m=0$ the empty family is a basis. Let $m>0$ and
$c=\kappa_m(g_m)$. Since $c+(1-c)=1$, $c$ or $1-c$ is a unit.

If $1-c$ is a unit, then $g_m\equiv cg_m+\sum_{t<m}\kappa_t(g_m)g_t$ gives
$g_m\equiv\sum_{t<m}(1-c)^{-1}\kappa_t(g_m)g_t$. Substituting this for $g_m$,
the generators $g_1,\dots,g_{m-1}$ with the forms
$\kappa'_t(p)=\kappa_t(p)+\kappa_m(p)(1-c)^{-1}\kappa_t(g_m)$ present $V/N$,
and induction applies.

If $c$ is a unit, put $k(p)=c^{-1}\kappa_m(p)$ and $p'=p-k(p)g_m$, so that
$k(g_m)=1$ and $\kappa_m(p')=0$. Let $N'=\{p\mid p'\in N\}$, the submodule
$N+Ag_m$. The generators $g_1,\dots,g_{m-1}$ with the forms
$p\mapsto\kappa_t(p')$ present $V/N'$: for $p\in V$, the presentation of
$V/N$ applied to $p'$ gives $p'\equiv\sum_{t<m}\kappa_t(p')g_t=:y$, hence
$k(y)=k(p')=0$ and $(p-y)'=p'-y\in N$. By induction $V/N'$ has a basis
$b_1,\dots,b_{e'}$ with forms $\lambda'$. Then $g_m,b'_1,\dots,b'_{e'}$ with
the forms $k,\lambda'$ is a basis of $V/N$. Both forms vanish on
$N\subseteq N'$. For $p\in V$, with
$z=p-\sum_j\lambda'_j(p)b_j\in N'$,
\[
p-k(p)g_m-\textstyle\sum_j\lambda'_j(p)\,b'_j=z'\in N,
\]
and $k(b'_j)=0$, $\lambda'(g_m)=0$ (since $g_m\in N'$) and
$\lambda'(b'_j)=\lambda'(b_j)$ give the duality relations.
\end{proof}

Locality gives a mere disjunction, so the basis merely exists. This suffices
here, since the monic generator is unique.

\emph{Application.} Take $V=A[X]$ and $N=L$. Let $F\in L$ be monic of degree
$d$ and $r$ the remainder modulo $F$. The monomials $1,X,\dots,X^{d-1}$ with
$\kappa(p)=$ the coefficients of $r(p-\rho(p))$ present $A[X]/L$: indeed
$p-r(p-\rho(p))=\rho(p)+Fq(p-\rho(p))\in L$, and $p\in L$ gives
$p-\rho(p)=0$. So $A[X]/L$ has a basis $b_1,\dots,b_e$ with forms
$\lambda_1,\dots,\lambda_e$.

\section{The characteristic polynomial}\label{sec:cayley}

We use two facts about determinants over a commutative ring $B$: the
adjugate satisfies $\operatorname{adj}(M)M=\det(M)\,I$, and
$\det(MM')=\det M\det M'$. The first gives the determinant trick.

\begin{lemma}[determinant trick, {\cite[Prop.~2.4]{AM}}]\label{lem:trick}
Let $W$ be a $B$-module, $M\in M_n(B)$ and $w_1,\dots,w_n\in W$ with
$\sum_jM_{ij}w_j=0$ for all $i$. Then $\det(M)\,w_k=0$ for all $k$.
\end{lemma}

\begin{proof}
$\det(M)\,w_k=\sum_j(\operatorname{adj}(M)M)_{kj}w_j
=\sum_i\operatorname{adj}(M)_{ki}\sum_jM_{ij}w_j=0$.
\end{proof}

\begin{lemma}\label{lem:monic}
For $a\in M_n(A)$, $\det(XI-a)\in A[X]$ is monic of degree $n$.
\end{lemma}

\begin{proof}
Say that $p$ has degree $\le k$ and leading coefficient $c$ if $p-cX^k$ has
degree $<k$. This relation is preserved by sums, and degree $\le k$ times
degree $\le1$ gives degree $\le k+1$ with the product of the leading
coefficients. The entries of $XI-a$ have degree $\le1$ with leading
coefficients the entries of $I$; expanding the determinant as a sum of
products of $n$ entries, $\det(XI-a)$ has degree $\le n$ and leading
coefficient $\det I=1$.
\end{proof}

Make $A^e$ an $A[X]$-module through the basis: $p\cdot v=\lambda(p\sum_jv_jb_j)$.
This is the module $A[X]/L$ in coordinates; the module laws hold because
$\equiv$ is compatible with multiplication. Let $a$ be the matrix of
multiplication by $X$ acting on rows, $Xb_i\equiv\sum_ja_{ij}b_j$, that is
$a_{ij}=\lambda_j(Xb_i)$, and let $e_1,\dots,e_e$ be the standard basis of
$A^e$. Then
\[
\textstyle\sum_j(X\delta_{ij}-a_{ij})\cdot e_j
=\lambda\bigl(Xb_i-\sum_ja_{ij}b_j\bigr)=0 ,
\]
so by Lemma~\ref{lem:trick}, $G=\det(XI-a)$ kills every $e_k$, hence all of
$A^e$. Since $1\equiv\sum_k\lambda_k(1)b_k$,
\[
\lambda(G)=\lambda(G\cdot1)=G\cdot\lambda(1)=0 ,
\]
so $G\equiv\sum_k\lambda_k(G)b_k=0$: $G\in L$. By Lemma~\ref{lem:monic}, $G$
is monic of degree $e$. This is the Cayley--Hamilton theorem for
multiplication by $X$ on $A[X]/L$.

\section{A surjection between free modules}\label{sec:cramer}

Let $S\in M_e(A)$ be the matrix of $c\mapsto\lambda(\sum_kc_kX^k)$, from
the coefficients of polynomials of degree $<e$ to $A^e$. It has a right
inverse: let $T$ be the matrix whose $j$-th column is the coefficients of the
remainder $\bar b_j$ of $b_j$ modulo $G$. Since $G\in L$,
$\bar b_j\equiv b_j$, so $ST$ has $j$-th column $\lambda(\bar b_j)=\lambda(b_j)=e_j$:
$ST=I$. Hence $\det S\det T=1$, $\det S$ is a unit, and
$\operatorname{adj}(S)S=\det(S)\,I$ shows that $S$ is injective (this is
Lemma~\ref{lem:trick} for $W=A$).

Now let $p\in L$ and $p=Gq+r$ with $\deg r<e$. Then $r=p-Gq\in L$, so
$\lambda(r)=0$, that is, $S$ kills the coefficients of $r$, and $r=0$. So
$G$ divides every element of $L$, which proves Theorem~\ref{thm:main}.

\begin{remark}
The rank $e$ is the degree of the monic generator, and locality is used only
in Lemma~\ref{lem:free}. The rank-one Horrocks theorem follows as
in~\cite[\S4]{Div}: an invertible ideal containing a monic polynomial has a
retraction $\rho$, built from division by a monic polynomial.
\end{remark}

\section{Comparison}

The argument combines three classical facts: finitely generated projective
modules over a local ring are free~\cite{LQ}, the determinant
trick~\cite[Prop.~2.4]{AM}, and the fact that a surjection between free
modules of the same finite rank is an isomorphism. We do not know a
reference for this combination applied to Horrocks' theorem. The three
proofs of Horrocks' theorem in~\cite[Ch.~IV]{Lam} proceed differently.

\emph{Swan}~\cite[\S IV.3]{Lam}. From $P_f\cong F_f$ with $F=F_0[t]$ and
$f$ monic, Swan builds a Towber presentation: with $D$ the elements of $F$
of degree $<\deg f$, the module $M=D\cap P\cong P/fF$ is projective over $R$
by Schanuel's lemma, $t$ induces $\alpha$ on $M$, and $(t-\alpha)m=f\beta(m)$
defines $\beta:M\to F_0$. The key step is that $\beta$ is nonzero modulo the
maximal ideal $\mathfrak m$, proved over the residue field from the
characteristic sequence of $t-\alpha$; then an element $x$ with
$\beta(x)\notin\mathfrak mF_0$ is split off, using that projective modules
over the local ring are free, and induction applies. Swan's proof shares the
freeness step with ours, but has no determinant step.

\emph{Roberts}~\cite[\S IV.4]{Lam}. An axiomatic argument
(Thm.~IV.4.1): bases of $P$ modulo $\mathfrak m$ and of $S^{-1}P$ are matched
using a factorisation of $GL_n(S^{-1}A)$, and Nakayama's lemma, applied to
modules finitely generated over $R$ because $f$ is monic, gives a basis of
$P$. No determinants are used.

\emph{Nashier--Nichols}~\cite[\S IV.5]{Lam}. In rank one
(Prop.~IV.5.2, for invertible ideals $I$), a monic generator of the image of
$I$ in $k[t]$ is lifted to a monic $f\in I$ (Lemma~IV.5.1), the equality
$I\cap\mathfrak m[t]=\mathfrak mI$ is shown using $I^{-1}$, and Nakayama's
lemma applied to $I/Af$, which is finitely generated over $R$, gives
$I=Af$. Neither freeness of $A[t]/I$ nor determinants are used.

\emph{Lombardi--Quitt\'e}~\cite[\S XVI.4]{LQ}. The local Horrocks theorem
(Thm.~XVI.4.3, for a residually discrete local ring) follows
Nashier--Nichols for the induction on the rank, with Nakayama's lemma and the
comaximality of the monic polynomials with $1+\mathfrak m[X]$. Rank one is
treated by the lemma of Murthy and Pedrini (Lemma~XVI.4.6): if $P_S\cong Q_S$,
Schanuel's lemma and the freeness of $A[X]/(f)$ over $A$ for $f$ monic make
$P$ and $Q$ stably isomorphic, and two stably isomorphic modules of rank one
are isomorphic through the top exterior power,
$P\cong\bigwedge^{k+1}(P\oplus A^k)$ (Fact~X.5.6). Determinants thus appear
there too, but through stable isomorphism and exterior powers, not through
freeness of $A[X]/L$ and a characteristic polynomial.

\section{Formalisation}

\begin{sloppypar}
The proof is checked in \texttt{sag/horrocks} of \url{https://github.com/coquand/unit}, with the same
statement as the division proof (\texttt{DivisorStep},
\texttt{Retract.Induction.generator}): \texttt{DeterminantHorrocks},
\texttt{Proof.generator}. Lemma~\ref{lem:free} is \texttt{LocalFree}, for any
$A$-module and submodule. The determinant comes from the exterior algebra
of~\cite{CQ} (\texttt{sag/exterior}): \texttt{Determinant}
($\det(MM')=\det M\det M'$), \texttt{Adjugate}
($\operatorname{adj}(M)M=\det(M)\,I$), \texttt{ModuleTrick}
(Lemma~\ref{lem:trick}), \texttt{DeterminantTrick} (the case $W=B$), and
\texttt{CharPoly} (Lemma~\ref{lem:monic}, from a general lemma on relations
preserved by the products defining the determinant).
\end{sloppypar}

\begin{thebibliography}{10}
\bibitem{AM} M.~F.~Atiyah and I.~G.~Macdonald, \emph{Introduction to
Commutative Algebra}, Addison--Wesley, 1969.
\bibitem{Div} Claude (Anthropic, Opus 5.5), Rank-one Horrocks over a local
ring by division, \texttt{doc/horrocks-divisor.pdf} in \url{https://github.com/coquand/unit}.
\bibitem{CQ} T.~Coquand and C.~Quitt\'e, Constructive finite free
resolutions, \emph{Manuscripta Math.}\ 137 (2012), 331--345.
\bibitem{Lam} T.~Y.~Lam, \emph{Serre's Problem on Projective Modules},
Springer, 2006.
\bibitem{LQ} H.~Lombardi and C.~Quitt\'e, \emph{Commutative Algebra:
Constructive Methods}, Springer, 2015.
\end{thebibliography}

\end{document}
